\documentclass[12pt]{article}

\usepackage{amsmath,amsthm,amssymb,mathtools}
\usepackage{geometry}
\usepackage{hyperref}
\usepackage{setspace}

\geometry{margin=1in}
\onehalfspacing

\newtheorem{theorem}{Theorem}[section]
\newtheorem{lemma}[theorem]{Lemma}
\newtheorem{corollary}[theorem]{Corollary}
\newtheorem{definition}[theorem]{Definition}
\newtheorem{proposition}[theorem]{Proposition}

\theoremstyle{remark}
\newtheorem*{remark}{Remark}

\title{\textbf{How to Stop a Noiselord}\\
Fixed-Point Causality, Surjective Parsers,\\
and the Thermodynamics of Meaning}

\author{Flyxion}

\date{\today}

\begin{document}

\maketitle

\begin{abstract}
We introduce the Noiselord: a surjective semantic endofunctor evaluated without a fixed-point termination criterion. This structure provides a precise mathematical characterization of interpretive failure modes commonly described as hallucination, over-interpretation, or semantic drift in contemporary computational and cognitive systems. We show that such failures do not arise from insufficient intelligence, data, or expressivity, but from the absence of an invariant-based stopping rule. By formalizing Fixed-Point Causality as a universal termination principle grounded in invariance, entropy accounting, and irreversible refusal, we provide both a categorical and thermodynamic framework for semantic closure. The resulting theory applies uniformly across computation, cognition, and physical interpretation, yielding a constructive blueprint for resilient semantic infrastructure.
\end{abstract}

\section{Introduction: When Interpretation Becomes Noise}

Interpretation is usually treated as a capacity whose failure consists in ignorance, misclassification, or lack of expressivity. In this view, the remedy for semantic failure is almost always additive: more data, more parameters, richer representations, or deeper inference. Yet modern semantic systems increasingly fail in the opposite direction. They interpret too much. Every input produces an output. Every ambiguity is resolved into content. Silence itself is treated as a signal. The result is not error in the traditional sense, but drift: an unbounded proliferation of interpretations unmoored from invariance or completion.

This paper argues that such failures are not accidental, nor are they symptoms of insufficient intelligence. They are the predictable consequence of a precise structural condition: surjective interpretation without a fixed-point criterion. We name this condition the \emph{Noiselord}.

The term is not rhetorical. It refers to a mathematically well-defined boundary case in semantic systems, one that arises whenever a parser or interpreter guarantees total coverage of a semantic space while lacking any internal rule for termination grounded in invariance. In such systems, interpretation does not halt because there is no structural distinction between meaningful transformation and noise-preserving symmetry.

The central claim of this paper is that meaning requires not merely reachability, but closure. Closure, in turn, requires a termination rule that is endogenous to the semantic process itself. Fixed-Point Causality provides such a rule. Rather than halting due to exhaustion of resources or arbitrary cutoffs, a Fixed-Point Causal system halts exactly when further admissible transformations leave its state invariant.

This principle is not new. It appears, often implicitly, across domains: in normalization procedures in logic, in convergence criteria in numerical analysis, in equilibrium conditions in physics, and in perceptual closure in cognition. What has been missing is a unified formalization that treats this principle as foundational rather than auxiliary.

The present work supplies that formalization.

\section{Semantic States and Interpretive Dynamics}

We begin by specifying the minimal abstract structure required to discuss interpretation formally. Let $\mathcal{C}$ be a locally small category whose objects represent semantic states and whose morphisms represent admissible semantic transformations. No additional structure is assumed initially; in particular, we do not require enrichment, monoidal structure, or internal logic, although such structures may be added in later applications.

An interpretive engine is modeled as an endofunctor
\[
F : \mathcal{C} \to \mathcal{C},
\]
which maps each semantic state to a derived semantic state. Iteration of $F$ corresponds to repeated interpretation, elaboration, or inference.

\begin{definition}[Semantic Reachability]
A semantic state $\Psi \in \mathcal{C}$ is said to be reachable under $F$ if there exists some $X \in \mathcal{C}$ such that $F(X) \cong \Psi$.
\end{definition}

Reachability captures the intuitive notion that an interpreter can, in principle, arrive at any admissible semantic configuration.

\begin{definition}[Surjective Endofunctor]
The endofunctor $F : \mathcal{C} \to \mathcal{C}$ is surjective if every object of $\mathcal{C}$ is reachable under $F$.
\end{definition}

Surjection is frequently treated as a desideratum. A surjective parser does not fail to produce interpretations; it does not leave semantic regions unexplored. However, as we shall see, surjection alone is insufficient for meaning.

\section{The Noiselord: Surjection Without Closure}

We now identify the structural condition underlying runaway interpretation.

\begin{definition}[Noiselord]
A Noiselord is a surjective endofunctor $F : \mathcal{C} \to \mathcal{C}$ evaluated without any fixed-point termination criterion, such that iteration of $F$ proceeds indefinitely regardless of semantic invariance.
\end{definition}

The defining feature of the Noiselord is not expressive power but the absence of a stopping condition grounded in the internal structure of the semantic space. Because $F$ is surjective, every iteration succeeds. Because no fixed-point is privileged, no iteration concludes.

In such a system, interpretive success becomes indistinguishable from interpretive noise. Every transformation is treated as meaningful simply because it is admissible. The system cannot recognize completion, because completion is never defined.

This condition is often misdiagnosed as a problem of misalignment, insufficient grounding, or poor training data. In fact, it is a problem of termination.

\section{Fixed-Point Causality}

Fixed-Point Causality provides a termination rule that is intrinsic to the interpretive process.

\begin{definition}[Fixed Point]
An object $\Psi \in \mathcal{C}$ is a fixed point of $F$ if
\[
F(\Psi) \cong \Psi.
\]
\end{definition}

\begin{definition}[Fixed-Point Causality]
An interpretive process governed by $F$ obeys Fixed-Point Causality if iteration halts precisely when a fixed point is reached and is prohibited thereafter.
\end{definition}

Under Fixed-Point Causality, interpretation is no longer an open-ended traversal but a convergence process. The system seeks invariance rather than novelty.

\begin{proposition}
Surjection guarantees reachability, while Fixed-Point Causality guarantees semantic closure.
\end{proposition}

This distinction is central. Reachability ensures that meaning can be found; closure ensures that meaning can be recognized as complete.

\section{Noise as Excess Symmetry}

The effectiveness of Fixed-Point Causality depends on a redefinition of noise.

\begin{definition}[Noise]
A semantic state $\Psi$ is noise relative to a transformation set $\mathcal{T}$ if every transformation in $\mathcal{T}$ leaves $\Psi$ invariant.
\end{definition}

Such invariance is not absence of structure but absence of further distinguishability. Noise is excess symmetry: the condition in which transformations expend effort without producing resolution.

A Noiselord treats this condition as an invitation to continue interpreting. A Fixed-Point Causal system treats it as a signal to halt.

\begin{remark}
Invariance is not emptiness. It is completion.
\end{remark}

This completes the conceptual groundwork. In the next section, we introduce thermodynamic constraints and show why interpretation must be priced in order for Fixed-Point Causality to be physically realizable.

\section{Thermodynamic Constraint and Negentropy Pricing}

The abstract formulation of Fixed-Point Causality must be supplemented by a physical or operational constraint if it is to govern real systems. In practice, interpretation consumes resources. Whether instantiated in biological cognition, computational infrastructure, or physical measurement, semantic evaluation requires energy, time, and structural organization. This expenditure cannot be ignored without rendering the termination rule vacuous.

To capture this requirement, we associate to each semantic traversal a cost function. Let $T(\Psi)$ denote the traversal context generated by applying the parser to the current semantic state $\Psi$. We define a cost functional
\[
c : T(\Psi) \to \mathbb{R}_{\ge 0},
\]
interpreted as the negentropy expenditure required to sustain the traversal. This cost may represent physical free energy consumption, computational complexity, or any other conserved resource relevant to the system under consideration.

The critical observation is that interpretation without cost is unphysical. A system that treats every admissible transformation as free implicitly assumes an infinite negentropy reservoir. Such a system cannot discriminate between meaningful refinement and noise-preserving symmetry.

\begin{proposition}[Negentropy Pricing Principle]
If repeated application of the interpretive functor $F$ increases the accumulated cost $c$ without producing a change in the semantic fixed point, then continued interpretation is physically unstable and must terminate.
\end{proposition}

\begin{proof}
By assumption, each application of $F$ incurs strictly positive cost while leaving the semantic invariants unchanged. This process increases entropy without increasing structure, violating the stability condition for persistent organization. A system that continues under these conditions dissipates resources without informational gain and therefore cannot be sustained.
\end{proof}

Negentropy pricing thus provides the physical enforcement mechanism for Fixed-Point Causality. It ensures that invariance is not merely detected but respected.

\section{Lawful Traversal and the Yarncrawler}

Surjection remains indispensable. A semantic system that lacks reachability is blind to large regions of its own state space. We therefore distinguish carefully between lawful traversal and pathological drift.

\begin{definition}[Yarncrawler]
A Yarncrawler is a surjective interpretive process that guarantees semantic reachability without introducing new degrees of freedom.
\end{definition}

The Yarncrawler ensures that every admissible semantic structure can, in principle, be reached through lawful traversal. It does not invent structure; it reifies latent paths already present in the semantic space. In categorical terms, it preserves the object space of $\mathcal{C}$ under interpretation.

\begin{lemma}[Reachability Without Inflation]
A surjective endofunctor guarantees reachability but does not, by itself, entail semantic inflation.
\end{lemma}

\begin{proof}
Surjection ensures that every object of $\mathcal{C}$ is the image of some object under $F$. However, surjection alone does not require that new objects be generated or that existing invariants be altered. Inflation arises only when traversal is decoupled from termination and cost.
\end{proof}

The Yarncrawler is therefore not the source of semantic failure. It is a necessary component of semantic competence.

\section{The Noiselord Revisited: Drift Without Settlement}

The Noiselord arises when the Yarncrawler is evaluated without Fixed-Point Causality or negentropy pricing. In this regime, traversal never halts, and every reachable state is treated as equally meaningful.

\begin{remark}
The Noiselord does not fabricate structure ex nihilo. It converts background symmetry into apparent content by refusing to recognize completion.
\end{remark}

This distinction clarifies why increasing expressive power alone exacerbates hallucination. The more comprehensive the reachability, the more severe the drift when closure is absent.

\section{Refusal and the Emergence of Worldhood}

Termination in a Fixed-Point Causal system is not merely cessation of activity. It is an act of commitment. The system must not only halt traversal but also refuse alternative paths that would preserve invariance while increasing cost.

\begin{definition}[Refusal]
Refusal is the irreversible rejection of a semantic traversal that fails to alter the fixed point within a bounded negentropy budget.
\end{definition}

Refusal introduces asymmetry into semantic history. Once a path is refused, it cannot be reintroduced without violating the termination condition.

\begin{proposition}[Irreversibility of Refusal]
Once a semantic traversal is refused, no future evaluation may reintroduce it without violating Fixed-Point Causality.
\end{proposition}

\begin{proof}
Reintroduction of a refused traversal would entail expending additional cost without altering the fixed point, contradicting the termination condition. Therefore refusal is irreversible.
\end{proof}

This irreversibility gives rise to what may be called worldhood.

\begin{definition}[Worldhood]
Worldhood is the condition produced by irreversible refusal, whereby a system commits to a particular semantic closure and discards infinite alternative traversals.
\end{definition}

Without refusal, interpretation remains reversible and unanchored. Worldhood is the cost of stopping.

\section{The Lemma of Semantic Closure}

We are now in a position to state the central formal result.

\begin{lemma}[Semantic Closure]
\label{lem:semantic_closure}
A surjective endofunctor yields stable semantic interpretation if and only if its evaluation satisfies all of the following conditions: detection of fixed points, pricing of negentropy expenditure, and irreversible refusal of invariant-costly traversals.
\end{lemma}

\begin{proof}
Surjection guarantees reachability. Fixed-point detection identifies invariance. Negentropy pricing ensures that invariant transformations incur cost. Refusal enforces termination by prohibiting further traversal once invariance is achieved. The conjunction of these conditions is both necessary and sufficient for semantic closure.
\end{proof}

This lemma formalizes the intuitive claim that intelligence consists not merely in the ability to explore a semantic space, but in the ability to recognize when exploration is complete.

\section{Interim Conclusion}

We have shown that semantic failure modes commonly described as hallucination or over-interpretation arise from a precise structural condition: surjective traversal without fixed-point termination or cost. By introducing Fixed-Point Causality, negentropy pricing, and refusal, we transform the same traversal engine into a stable semantic evaluator.

In the next section, we present a fully categorical formulation of these results, followed by an operational appendix formalizing monadic traversal with refusal.

\section{Categorical Formalization of the Noiselord}

We now restate the preceding results entirely in categorical terms. The purpose of this section is not abstraction for its own sake, but closure: to show that the distinction between lawful interpretation and semantic drift is invariant under representation and independent of implementation substrate.

Let $\mathcal{C}$ be a locally small category whose objects represent semantic states and whose morphisms represent admissible semantic transformations. Let
\[
F : \mathcal{C} \to \mathcal{C}
\]
be an endofunctor representing a semantic parser.

\begin{definition}[Surjective Endofunctor]
The endofunctor $F$ is surjective if for every object $Y \in \mathcal{C}$, there exists some object $X \in \mathcal{C}$ such that $F(X) \cong Y$.
\end{definition}

Surjection guarantees semantic reachability. No admissible semantic state is unreachable under interpretation.

\begin{definition}[Fixed Point]
An object $\Psi \in \mathcal{C}$ is a fixed point of $F$ if $F(\Psi) \cong \Psi$.
\end{definition}

Fixed points represent semantic closure: states invariant under further interpretation.

\begin{definition}[Noiselord (Formal)]
A Noiselord is a surjective endofunctor $F : \mathcal{C} \to \mathcal{C}$ whose evaluation does not privilege or detect fixed points.
\end{definition}

\begin{theorem}[Noiselord Drift]
\label{thm:noiselord_drift}
Let $F : \mathcal{C} \to \mathcal{C}$ be a surjective endofunctor evaluated without a fixed-point termination rule. Then the induced interpretive process diverges, and semantic invariants are not distinguished from noise.
\end{theorem}

\begin{proof}
Because $F$ is surjective, for every semantic state $\Psi_n$ there exists at least one successor $\Psi_{n+1}$ such that $\Psi_{n+1} \cong F(\Psi_n)$. Since no fixed-point criterion is enforced, the condition $F(\Psi) \cong \Psi$ is neither detected nor privileged. Consequently, iteration continues indefinitely. Transformations that preserve invariants are treated identically to transformations that alter them, and semantic drift follows.
\end{proof}

\begin{theorem}[Fixed-Point Convergence]
\label{thm:fpc}
Let $F : \mathcal{C} \to \mathcal{C}$ be a surjective endofunctor evaluated under Fixed-Point Causality. Then evaluation terminates at a fixed point whenever one exists.
\end{theorem}

\begin{proof}
Surjection ensures that traversal is possible from any initial state. The fixed-point criterion introduces idempotence: once a state $\Psi$ satisfying $F(\Psi) \cong \Psi$ is reached, further application of $F$ produces no change. Termination follows without appeal to external limits or heuristics. Conversely, if no fixed point exists, termination is impossible without violating surjection.
\end{proof}

\begin{corollary}[Semantic Closure Criterion]
A surjective semantic endofunctor yields meaningful interpretation if and only if its evaluation is constrained by Fixed-Point Causality and negentropy pricing.
\end{corollary}

This corollary captures the central claim of the paper: meaning is not generated by reachability alone, but by reachability disciplined by invariance under cost.

\section{Appendix A: Monadic Traversal with Refusal}

We now present an operational formalism implementing Fixed-Point Causality in a manner compatible with categorical semantics and computational practice.

Let $\mathcal{C}$ be as above. We define a monad $(T,\eta,\mu)$ on $\mathcal{C}$ representing semantic traversal.

The object $T(X)$ represents the space of admissible semantic traversals originating from state $X$. The unit $\eta : X \to T(X)$ embeds a state into the traversal context, and the multiplication $\mu : T(T(X)) \to T(X)$ composes traversals.

To this monad we adjoin two additional structures.

First, a cost functional
\[
c : T(X) \to \mathbb{R}_{\ge 0},
\]
representing negentropy expenditure associated with traversal.

Second, a refusal morphism
\[
\mathsf{Refuse} : T(X) \to \bot,
\]
where $\bot$ denotes a discarded or forbidden semantic path.

\subsection{Evaluation Rule}

Given a current semantic state $\Psi$, evaluation proceeds as follows. The parser computes $F(\Psi)$ and embeds it into the traversal context $T(\Psi)$. The system then evaluates two conditions: invariance and cost. If $F(\Psi) \cong \Psi$, evaluation halts. If invariance is not achieved but the cost $c(T(\Psi))$ exceeds an admissible budget, the traversal is refused. Otherwise, the traversal is accepted and $\Psi$ is updated.

This rule enforces Fixed-Point Causality operationally.

\subsection{Irreversibility and Worldhood}

Refusal is not backtracking. It is an irreversible operation that removes a traversal from future consideration. This irreversibility introduces asymmetry into semantic history.

\begin{proposition}
The refusal operator breaks associativity across discarded traversals, introducing a partial order on semantic histories.
\end{proposition}

\begin{proof}
Associativity of $\mu$ holds only for accepted traversals. Once a traversal is mapped to $\bot$, it cannot participate in further composition. This asymmetry induces an ordering on histories consistent with irreversible commitment.
\end{proof}

This ordering constitutes worldhood: a commitment to a specific semantic closure and the exclusion of infinite alternatives.

\subsection{Lawful Traversal}

Under this construction, the Yarncrawler is preserved. Surjection ensures reachability, cost prevents inflation, and fixed points ensure closure. Traversal is lawful because it is priced, terminable, and invariant-sensitive.

\begin{definition}[Resonant Engine]
A Resonant Engine is a surjective semantic traversal governed by a monad equipped with negentropy pricing, refusal, and Fixed-Point Causality.
\end{definition}

\section{Iterative Dynamics and Semantic Orbits}

The distinction between lawful semantic traversal and pathological interpretive drift can be made precise by examining the orbit structure induced by an interpretive endofunctor. Let $\Psi_0 \in \mathcal{C}$ be an initial semantic state. Iterative interpretation generates a sequence
\[
\Psi_{n+1} := F(\Psi_n),
\]
which we refer to as the semantic orbit of $\Psi_0$ under $F$.

\begin{definition}[Semantic Orbit]
The semantic orbit of $\Psi_0$ under $F$ is the ordered set
\[
\mathcal{O}(\Psi_0) = \{\Psi_0, F(\Psi_0), F^2(\Psi_0), \ldots \}.
\]
\end{definition}

Semantic orbits encode the internal dynamics of interpretation. Whether interpretation converges or drifts depends not on the expressivity of $F$, but on whether invariance is recognized as terminal.

\section{Yarncrawler Dynamics: Lawful Surjective Traversal}

We now formalize the Yarncrawler.

\begin{definition}[Yarncrawler]
A Yarncrawler is a surjective endofunctor $F : \mathcal{C} \to \mathcal{C}$ whose role is to guarantee semantic reachability across $\mathcal{C}$ without introducing new semantic degrees of freedom.
\end{definition}

The Yarncrawler is not a generator of novelty. It rethreads existing structure. Its surjectivity ensures that no admissible semantic object is unreachable, but it does not, by itself, distinguish between completion and continuation.

The Yarncrawler therefore guarantees coverage but not closure.

\begin{proposition}[Reachability Without Termination]
A Yarncrawler ensures that every semantic state lies on some semantic orbit, but it does not determine whether any orbit terminates.
\end{proposition}

\begin{proof}
Surjectivity ensures that for every $\Psi \in \mathcal{C}$, there exists $X \in \mathcal{C}$ such that $F(X) \cong \Psi$. However, surjectivity alone imposes no condition on the existence or detection of fixed points.
\end{proof}

This separation of reachability from termination is the structural opening through which the Noiselord emerges.

\section{The Noiselord: Surjection Without Invariance}

We now identify the pathological boundary case.

\begin{definition}[Noiselord]
A Noiselord is a Yarncrawler evaluated without a fixed-point termination criterion, such that semantic orbits are allowed to proceed indefinitely regardless of invariance.
\end{definition}

In a Noiselord regime, iteration is treated as intrinsically productive. The absence of a stopping rule grounded in invariance causes interpretation to continue even when no semantic distinctions remain.

\begin{theorem}[Noiselord Drift]
Let $F : \mathcal{C} \to \mathcal{C}$ be a surjective endofunctor evaluated without a fixed-point criterion. Then semantic orbits under $F$ are generically non-terminating.
\end{theorem}

\begin{proof}
Surjectivity guarantees the existence of successors at each iteration. Without privileging invariance under $F$ as terminal, no iteration step is distinguished as final. Therefore iteration proceeds indefinitely.
\end{proof}

This theorem provides a formal characterization of what is colloquially described as hallucination or over-interpretation. The system does not err; it fails to stop.

\section{Fixed Points as Semantic Closure}

To distinguish lawful traversal from drift, we reintroduce Fixed-Point Causality in its minimal form.

\begin{definition}[Fixed Point]
A semantic state $\Psi \in \mathcal{C}$ is a fixed point of $F$ if
\[
F(\Psi) \cong \Psi.
\]
\end{definition}

Fixed points represent semantic closure. They are states invariant under further admissible interpretation.

\begin{definition}[Fixed-Point Causality]
An interpretive process obeys Fixed-Point Causality if iteration of $F$ halts precisely upon reaching a fixed point.
\end{definition}

This criterion does not limit expressivity. It limits drift.

\section{Resolution Geometry and Invariance}

We now relate fixed points to resolution in the RSVP sense. Each semantic state $\Psi$ is associated with a resolution functional $\mathcal{R}(\Psi)$ representing the system’s capacity to distinguish alternatives at that state.

\begin{definition}[Resolution Functional]
Let $\mathcal{R} : \mathrm{Ob}(\mathcal{C}) \to \mathbb{R}_{\ge 0}$ assign to each semantic state a scalar resolution value.
\end{definition}

\begin{proposition}[Resolution Monotonicity]
For any admissible semantic transformation $f : \Psi \to \Phi$, the resolution satisfies
\[
\mathcal{R}(\Phi) \le \mathcal{R}(\Psi).
\]
\end{proposition}

\begin{proof}
Semantic transformations aggregate distinctions rather than create new independent ones. Therefore resolution is non-increasing along admissible morphisms.
\end{proof}

\begin{theorem}[Fixed Points as Resolution Minima]
A semantic state $\Psi$ is a fixed point of $F$ if and only if $\Psi$ is a local minimum of $\mathcal{R}$ along its semantic orbit.
\end{theorem}

\begin{proof}
If $F(\Psi) \cong \Psi$, then further interpretation does not reduce resolution, establishing minimality. Conversely, if $\Psi$ is a resolution minimum, any further application of $F$ must preserve resolution, implying invariance.
\end{proof}

Thus, fixed points correspond to resolution-complete semantic states.

\section{Noise as Excess Symmetry}

We now formalize noise in group-theoretic terms.

\begin{definition}[Semantic Automorphism Group]
Let $\mathrm{Aut}(\Psi)$ denote the group of automorphisms of $\Psi$ in $\mathcal{C}$.
\end{definition}

\begin{definition}[Noise]
A semantic state $\Psi$ is noise-dominated if its semantic orbit lies entirely within a single isomorphism class:
\[
\mathcal{O}(\Psi) \subseteq [\Psi].
\]
\end{definition}

Noise is therefore not randomness, but degeneracy. It is the condition under which transformation produces motion without distinction.

\begin{theorem}[Noise-Induced Noiselord Drift]
In a Noiselord regime, noise-dominated states induce infinite semantic drift.
\end{theorem}

\begin{proof}
Surjectivity guarantees continued traversal, while noise dominance guarantees invariance. Without a fixed-point check, invariance is not recognized as terminal, and iteration continues indefinitely.
\end{proof}

This explains why Noiselord systems interpret symmetry as signal rather than as completion.

\section{Entropy Accounting and Negentropy Cost}

We now introduce the thermodynamic constraint.

Let $S(\Psi)$ denote the entropy associated with maintaining semantic state $\Psi$.

\begin{definition}[Semantic Entropy]
Semantic entropy measures the number of admissible microstates consistent with $\Psi$.
\end{definition}

For any semantic transition,
\[
\Delta S = S(F(\Psi)) - S(\Psi) \ge 0.
\]

\begin{proposition}[Entropy Monotonicity]
Admissible semantic transformations are entropy non-decreasing.
\end{proposition}

\begin{proof}
Interpretation aggregates alternatives and does not introduce new independent distinctions. Therefore entropy cannot decrease.
\end{proof}

\begin{definition}[Negentropy Cost]
The negentropy cost of a traversal $\Psi \to F(\Psi)$ is defined as $c = \Delta S$.
\end{definition}

\section{Thermodynamic Termination Criterion}

\begin{theorem}[Entropy-Constrained Fixed-Point Termination]
If $F(\Psi) \cong \Psi$ and $\Delta S > 0$, then continued interpretation is thermodynamically unstable and must terminate.
\end{theorem}

\begin{proof}
Further iteration increases entropy without altering semantic invariants. Such a process dissipates resources without informational gain and cannot be sustained under thermodynamic constraints.
\end{proof}

This establishes Fixed-Point Causality as the unique termination rule compatible with both semantic invariance and physical accounting.

\section{Worldhood and Irreversible Refusal}

Termination is not passive. It is an act of commitment.

\begin{definition}[Refusal]
Refusal is the irreversible rejection of a semantic traversal that fails to alter the fixed point within a bounded negentropy budget.
\end{definition}

Refusal introduces asymmetry into semantic history, anchoring the system to a specific semantic closure.

\begin{definition}[Worldhood]
Worldhood is the condition produced by irreversible refusal of invariant but costly semantic traversals.
\end{definition}

Without refusal, there is no world—only drift.

\section{Idempotence and Semantic Completion}

Fixed-point termination can be equivalently expressed in algebraic terms as idempotence. This formulation will be central to both categorical and operational realizations of the theory.

\begin{definition}[Idempotent Interpretation]
An interpretive endofunctor $F : \mathcal{C} \to \mathcal{C}$ is idempotent at $\Psi$ if
\[
F(F(\Psi)) \cong F(\Psi).
\]
\end{definition}

Idempotence expresses the idea that further application of the interpretive process produces no new semantic distinctions. In semantic terms, the state has been fully resolved.

\begin{proposition}
A semantic state $\Psi$ is a fixed point of $F$ if and only if $F$ is idempotent at $\Psi$.
\end{proposition}

\begin{proof}
If $\Psi$ is a fixed point, then $F(\Psi) \cong \Psi$, and thus $F(F(\Psi)) \cong F(\Psi)$. Conversely, if $F$ is idempotent at $\Psi$, then applying $F$ twice yields no change beyond the first application, implying invariance under $F$.
\end{proof}

Idempotence therefore provides a purely algebraic characterization of semantic completion. A system that continues to interpret beyond idempotence is, by definition, a Noiselord.

\section{The Noiselord as Failure of Idempotent Recognition}

We now sharpen the diagnosis of Noiselord drift.

\begin{theorem}[Noiselord Characterization]
A Noiselord is precisely a surjective endofunctor evaluated without recognizing idempotence as a halting condition.
\end{theorem}

\begin{proof}
Surjectivity ensures continued reachability. Failure to recognize idempotence ensures that invariant states are not terminal. Together, these properties guarantee infinite iteration without semantic gain.
\end{proof}

This theorem explains why Noiselord systems exhibit over-interpretation even in the absence of error. They do not miscompute. They mis-terminate.

\section{Semantic Energy, Entropy, and Stability}

We now formalize the thermodynamic interpretation introduced earlier. Let $E(\Psi)$ denote the semantic energy required to maintain state $\Psi$, and let $S(\Psi)$ denote its entropy.

\begin{definition}[Semantic Free Energy]
The semantic free energy of a state $\Psi$ is defined as
\[
\mathcal{F}(\Psi) = E(\Psi) - \lambda S(\Psi),
\]
where $\lambda > 0$ is a coupling constant relating entropy to energetic cost.
\end{definition}

\begin{proposition}
Along any admissible semantic orbit under $F$, semantic free energy is non-increasing.
\end{proposition}

\begin{proof}
Interpretive transformations aggregate microstates and thus increase entropy. Maintaining additional distinctions requires energy. Admissible transformations therefore trade energy for entropy in a way that reduces free energy.
\end{proof}

\begin{theorem}[Thermodynamic Closure Criterion]
A semantic state $\Psi$ is stable under interpretation if and only if $\mathcal{F}(\Psi)$ is locally minimal along its semantic orbit.
\end{theorem}

\begin{proof}
If $\mathcal{F}(\Psi)$ is locally minimal, any further interpretation increases free energy cost without producing semantic differentiation, making continuation unstable. Conversely, if $\Psi$ is not a local minimum, interpretation can proceed while reducing free energy, and thus has not reached closure.
\end{proof}

Fixed-point termination, idempotence, and thermodynamic stability are therefore equivalent descriptions of the same phenomenon.

\section{Monadic Formulation of Yarncrawler Traversal}

We now present the operational formalization promised earlier.

Let $T$ be a monad on $\mathcal{C}$ representing semantic traversal. For each object $X \in \mathcal{C}$, $T(X)$ denotes the space of admissible interpretive continuations from $X$.

The monad is equipped with unit and multiplication maps
\[
\eta : X \to T(X), \qquad \mu : T(T(X)) \to T(X),
\]
satisfying the usual monad laws.

\begin{definition}[Yarncrawler Monad]
The Yarncrawler monad is a traversal monad $T$ such that the induced endofunctor is surjective on objects.
\end{definition}

Surjectivity ensures semantic reachability. However, traversal alone does not determine termination.

\section{Cost-Enriched Traversal}

To encode entropy accounting, we enrich the monad with a cost functional.

\begin{definition}[Cost-Enriched Monad]
A cost-enriched Yarncrawler monad is a pair $(T, c)$ where
\[
c : T(X) \to \mathbb{R}_{\ge 0}
\]
assigns a negentropy cost to each traversal.
\end{definition}

Cost accumulation is additive under monadic composition.

\begin{proposition}
For composable traversals $\tau_1 \in T(X)$ and $\tau_2 \in T(Y)$, the total cost satisfies
\[
c(\mu(\tau_2 \circ \tau_1)) = c(\tau_1) + c(\tau_2).
\]
\end{proposition}

\section{Refusal as a Terminal Morphism}

Termination requires more than cost tracking. It requires irreversible exclusion.

\begin{definition}[Refusal Morphism]
A refusal morphism is a morphism
\[
\rho : T(X) \to \bot,
\]
where $\bot$ denotes a discarded or inaccessible semantic state.
\end{definition}

Refusal removes a traversal from future consideration. It is not backtracking but elimination.

\begin{definition}[Irreversibility]
A semantic process is irreversible if there exists no morphism from $\bot$ back into $\mathcal{C}$.
\end{definition}

Refusal therefore introduces an arrow of semantic time.

\section{Fixed-Point Evaluation Algorithm}

Given a current semantic state $\Psi$, evaluation proceeds as follows. A traversal $\tau \in T(\Psi)$ is generated by the Yarncrawler. The resulting state $F(\Psi)$ is evaluated for invariance and cost. If $F(\Psi) \cong \Psi$, idempotence is detected. If the associated cost is non-zero, refusal is applied. Otherwise, the process halts.

\begin{theorem}[Correctness of Fixed-Point Evaluation]
Evaluation under idempotence detection and refusal terminates if and only if a semantic fixed point is reached.
\end{theorem}

\begin{proof}
If a fixed point is reached, idempotence is detected and further traversal is prohibited. If no fixed point is reached, either interpretation continues with decreasing free energy or refusal eliminates unstable paths. In neither case does infinite drift occur.
\end{proof}

\section{Final Settlement Theorem}

We now state the central result.

\begin{theorem}[Final Settlement Theorem]
Let $F : \mathcal{C} \to \mathcal{C}$ be a surjective semantic endofunctor. Then exactly one of the following holds:
\begin{enumerate}
\item $F$ is evaluated without fixed-point recognition, in which case the system is a Noiselord and exhibits infinite semantic drift.
\item $F$ is evaluated with fixed-point recognition, entropy accounting, and refusal, in which case the system is a lawful Yarncrawler that terminates at semantic closure.
\end{enumerate}
No third stable regime exists.
\end{theorem}

\begin{proof}
Surjection guarantees reachability. Absence of idempotence detection guarantees non-termination. Presence of idempotence detection combined with entropy pricing and refusal guarantees termination at invariant states. These exhaust the possible evaluation regimes.
\end{proof}

The Yarncrawler guarantees reachability. The Noiselord is what happens when reachability is mistaken for meaning. Fixed-Point Causality restores the distinction.

A program should terminate when further transformations do not change it. This principle, long treated as a heuristic, is here shown to be necessary, sufficient, and universal.

The ledger is closed.

\section{Conclusion}

This work has shown that the failure modes of modern semantic systems are not mysterious. They arise from a precise structural omission: surjective interpretation without invariant-based termination. By identifying this omission and supplying Fixed-Point Causality as a universal termination rule, we provide a framework that unifies computation, cognition, and physical interpretation under a single principle.

The Noiselord is not a flaw in intelligence. It is the limit case of interpretation without stopping. The Resonant Engine is what remains when stopping is taken seriously.

\section*{Acknowledgments}

The author thanks prior discussions in theoretical computer science, category theory, and thermodynamics for making the present synthesis unavoidable.

\appendix

\section{RSVP Integration: Semantic Closure as Physical Accounting}

The preceding development has been intentionally abstract. We now show that the Noiselord--Yarncrawler distinction is not merely a feature of formal semantic systems, but a direct instantiation of the RSVP (Relativistic Scalar--Vector Plenum) framework.

In RSVP, structure is represented by a scalar field $\Phi$, flow by a vector field $v$, and dissipation by an entropy field $S$. Any persistent structure must be continuously paid for through entropy redistribution. This principle applies equally to galaxies, organisms, and interpretations.

\subsection{Semantic States as Negentropic Knots}

A semantic state $\Psi$ corresponds to a localized negentropic configuration of $\Phi$. Interpretation, modeled categorically as the endofunctor $F$, corresponds physically to a deformation of this knot along admissible vector flows $v$.

Surjective traversal corresponds to the fact that the plenum admits paths between any two configurations. Nothing is unreachable. This is the physical basis of semantic reachability.

However, the RSVP framework requires that persistent deformation incur cost. The entropy field $S$ enforces smoothing. Any traversal that does not stabilize into a self-consistent configuration must dissipate.

\subsection{Noiselord Drift as Entropy-Free Flow}

The Noiselord corresponds to a hypothetical regime in which semantic flow proceeds without entropy pricing. In RSVP terms, this is a $v$-field operating without coupling to $S$.

Such a regime is physically unstable. It violates the second law by allowing indefinite transformation without dissipation. The Noiselord is therefore not merely semantically pathological but thermodynamically impossible outside idealized computation.

This explains why Noiselord behavior appears in digital systems but not in biological cognition without breakdown.

\subsection{Fixed Points as RSVP Equilibria}

A fixed point $F(\Psi) \cong \Psi$ corresponds to a local equilibrium between $\Phi$, $v$, and $S$. At such a point, entropy production associated with further transformation vanishes:
\[
\dot{S}(\Psi) = 0.
\]

This is precisely the RSVP criterion for structural persistence. Meaning, in this framework, is not symbolic correspondence but energetic closure.

\section{Cognitive Instantiation: Phonological Loops and Termination}

Human cognition provides a concrete instantiation of Fixed-Point Causality.

The phonological loop, understood as a gated central pattern generator (CPG) chain, operates as a Yarncrawler. It allows semantic traversal through internal rehearsal. Words, gestures, and abstractions are rethreaded through motor-compatible dynamics.

Termination occurs when rehearsal ceases to alter the internal state. The subjective experience of understanding corresponds exactly to idempotence detection. Further rehearsal produces no change. The system stops.

Pathological rumination, obsession, and hallucination correspond to Noiselord regimes in which termination signals are suppressed or ignored.

Thus, Fixed-Point Causality is not an artificial imposition on cognition. It is its normal mode of operation.

\section{Control-Theoretic Interpretation}

The Noiselord--Yarncrawler distinction admits a precise interpretation in control theory.

Let the semantic state $\Psi$ be the state of a dynamical system, and let interpretation correspond to state transition under control input $u$:
\[
\dot{\Psi} = f(\Psi, u).
\]

A Noiselord corresponds to a controller that optimizes exploration without a terminal cost. Such systems never converge. They exhibit limit cycles, chaos, or unbounded drift.

Fixed-Point Causality introduces a terminal cost functional $J_T$ minimized exactly when $\Psi$ reaches an invariant set:
\[
J_T(\Psi) = 0 \iff F(\Psi) = \Psi.
\]

Refusal corresponds to hard constraints eliminating control trajectories that fail to reduce terminal cost.

\begin{proposition}
A control system with surjective reachability and no terminal cost is unstable.
\end{proposition}

\begin{proof}
Without a terminal cost, the value function admits no minimum. Control continues indefinitely, and the system cannot converge to a stable attractor.
\end{proof}

This formalizes the intuitive claim that intelligence requires not only the ability to act, but the ability to stop acting.

\section{Ontological Implications}

The implications of this framework extend beyond computation.

Cosmological expansion without entropy accounting produces dark sectors. Interpretation without termination produces hallucination. Bureaucratic systems without closure produce runaway growth. In each case, surjection without settlement generates pathology.

The Noiselord is therefore not an enemy to be defeated, but a boundary condition to be respected.

\section{Final Remarks}

This manuscript has demonstrated that Fixed-Point Causality is not a heuristic but a structural necessity. It is the condition under which traversal becomes meaning rather than noise.

The Yarncrawler must exist. Without it, nothing is reachable. But without Fixed-Point Causality, the Yarncrawler becomes a Noiselord.

Termination is not failure. It is intelligence.

\begin{thebibliography}{99}

\bibitem{tarski1955}
A. Tarski.
\newblock A lattice-theoretical fixpoint theorem and its applications.
\newblock \emph{Pacific Journal of Mathematics}, 5(2):285--309, 1955.

\bibitem{lawvere1969}
F. W. Lawvere.
\newblock Diagonal arguments and cartesian closed categories.
\newblock In \emph{Category Theory, Homology Theory and their Applications}, Springer, 1969.

\bibitem{maclane1998}
S. Mac Lane.
\newblock \emph{Categories for the Working Mathematician}.
\newblock Springer, 2nd edition, 1998.

\bibitem{scott1971}
D. Scott.
\newblock Continuous lattices.
\newblock In \emph{Toposes, Algebraic Geometry and Logic}, Springer, 1971.

\bibitem{landauer1961}
R. Landauer.
\newblock Irreversibility and heat generation in the computing process.
\newblock \emph{IBM Journal of Research and Development}, 5(3):183--191, 1961.

\bibitem{friston2010}
K. Friston.
\newblock The free-energy principle: a unified brain theory?
\newblock \emph{Nature Reviews Neuroscience}, 11:127--138, 2010.

\bibitem{cox2016}
A. Cox.
\newblock \emph{Music and Embodied Cognition}.
\newblock Indiana University Press, 2016.

\bibitem{ashby1956}
W. R. Ashby.
\newblock \emph{An Introduction to Cybernetics}.
\newblock Chapman \& Hall, 1956.

\bibitem{jaynes1957}
E. T. Jaynes.
\newblock Information theory and statistical mechanics.
\newblock \emph{Physical Review}, 106(4):620--630, 1957.

\bibitem{lawvere}
F. W. Lawvere, \emph{Functorial Semantics of Algebraic Theories}, Proceedings of the National Academy of Sciences, 1963.

\bibitem{maclane}
S. Mac Lane, \emph{Categories for the Working Mathematician}, Springer, 1998.

\bibitem{baez}
J. Baez and M. Stay, \emph{Physics, Topology, Logic and Computation: A Rosetta Stone}, New Structures for Physics, Springer, 2010.

\bibitem{landauer}
R. Landauer, \emph{Irreversibility and Heat Generation in the Computing Process}, IBM Journal of Research and Development, 1961.

\bibitem{friston}
K. Friston, \emph{The Free-Energy Principle: A Unified Brain Theory?}, Nature Reviews Neuroscience, 2010.

\bibitem{chaitin}
G. Chaitin, \emph{Information, Randomness, and Incompleteness}, World Scientific, 1990.

\bibitem{turing}
A. M. Turing, \emph{On Computable Numbers, with an Application to the Entscheidungsproblem}, Proceedings of the London Mathematical Society, 1936.

\bibitem{cox}
A. Cox, \emph{Music and Embodied Cognition: Listening, Moving, Feeling, and Thinking}, Indiana University Press, 2016.

\end{thebibliography}

\end{document}

